\section{Primitive packet and the dynamic residual support}
\label{sec:interface}

Fix \(h\in\Z\setminus\{0\}\), put \(H=\rad|h|\), and choose
\begin{equation}
 R=\lfloor X^{1/2-\delta}\rfloor,
 \qquad 0<\delta<\frac12.
 \label{eq:R}
\end{equation}
The finite Ramanujan model and residual are
\begin{equation}
 \Lambda_R(n)=\sum_{q\le R}\frac{\mu(q)}{\varphi(q)}c_q(n),
 \qquad r_R(n)=\Lambda(n)-\Lambda_R(n).
 \label{eq:finite-model}
\end{equation}
In one separated primitive row packet,
\begin{equation}
 Q=LD,\qquad QJ\asymp X,\qquad
 L>\max\{2R,2|h|\},
 \label{eq:row-scales}
\end{equation}
and the rows are
\begin{equation}
 m=\ell d\asymp Q,\qquad
 \ell\in[L/2,2L]\text{ prime},\qquad
 d\in[D,2D],\quad d\le R^{1/2}.
 \label{eq:rows}
\end{equation}
We impose \((m,H)=1\), \((j,H)=1\), and put
\begin{equation}
 n_m(j)=mj+h.
 \label{eq:target}
\end{equation}
Then
\begin{equation}
 (n_m(j),mH)=1.
 \label{eq:target-coprime}
\end{equation}
In particular, every divisor of a target is coprime to the appropriate
row and to \(H\); this exact fact replaces any unnecessary comparison
between that divisor and the source prime \(\ell\).

Let \(U\asymp X\) be a finite horizon containing all positive targets
in the separated packet.  Define
\begin{align}
 \lambda'_R(u)
 &=u\sum_{b\le R/u}\frac{\mu(ub)\mu(b)}{\varphi(ub)}
 \quad(u\le R),
 &\lambda'_R(u)&=0\quad(u>R),
 \label{eq:lambda-prime}\\
 a(u)&=-\mu(u)\log u,
 &b_R(u)&=a(u)-\lambda'_R(u).
 \label{eq:a-b}
\end{align}
Both coefficients are zero off squarefree integers, and for every
fixed \(\eps>0\), uniformly for \(u\le U\),
\begin{equation}
 |b_R(u)|\ll_\eps X^\eps.
 \label{eq:b-soft}
\end{equation}
Indeed, for \(u>R\) this is the elementary logarithmic bound, while for
\(u\le R\) it follows from the positive divisor formula for
\(\lambda'_R\).

TPC-20 proves the pointwise matched representation
\begin{equation}
 Z_m(j)=\eps_U(m)+
 \sum_{\substack{u\le U\\(u,mH)=1}}
 b_R(u)\left(\one_{u\mid n_m(j)}-\frac1u\right),
 \label{eq:matched-signal}
\end{equation}
where, for every \(A>0\),
\begin{equation}
 \eps_U(m)\ll_{A,h}X^\eps(\log X)^{-A}.
 \label{eq:calibration-defect}
\end{equation}
No prime number theorem in a growing row progression is used in this
calibration.

Choose a dynamic residual cutoff
\begin{equation}
 R<S\le U,
 \qquad S=X^{s+o(1)},
 \label{eq:S}
\end{equation}
and put
\begin{align}
 \vartheta_{R,S}(u)&=b_R(u)\one_{u\le S},
 \label{eq:theta-RS}\\
 Z_m^{\le S}(j)
 &=\sum_{\substack{u\le S\\(u,mH)=1}}
 b_R(u)\left(\one_{u\mid n_m(j)}-\frac1u\right),
 \label{eq:truncated-signal}\\
 Z_m^{>S}(j)
 &=\sum_{\substack{S<u\le U\\(u,mH)=1}}
 a(u)\left(\one_{u\mid n_m(j)}-\frac1u\right).
 \label{eq:ultra-signal}
\end{align}
Then \(Z_m=\eps_U(m)+Z_m^{\le S}+Z_m^{>S}\) exactly.  Notice that
\(S\) truncates the residual expansion only; the model
\eqref{eq:finite-model} still ends at \(R\).

The separated row weights \(x_m,y_n\) satisfy
\begin{equation}
 \|x\|_\infty+\|y\|_\infty\ll_\eps X^\eps,
 \qquad
 \sum_m|x_m|^2+\sum_n|y_n|^2\ll QX^\eps.
 \label{eq:row-energy}
\end{equation}
The actual generic Schur mask \(\mathfrak m(m,n)\) is independent of
all divisor variables, obeys
\begin{equation}
 |\mathfrak m(m,n)|\le1,
 \qquad \mathfrak m(m,m)=0,
 \label{eq:mask}
\end{equation}
and may include the distinct-source, determinant-separation, and
row-gcd restrictions of TPC-18.  The dynamic theorem below applies to
this row-pair-only mask.  A modulus-dependent mask would require a new
argument.

For squarefree \(u,v\le S\), write
\begin{equation}
 g=(u,v),\qquad u=ga,\qquad v=gb,\qquad
 q=[u,v]=gab.
 \label{eq:gab}
\end{equation}
The compatible joint divisibility class exists precisely when
\(g\mid m-n\), in addition to the primitive coprimality conditions.
To make the dynamic scalar object explicit, put
\(\cU_q=(\Z/q\Z)^\times\) and, for \(\xi\in\cU_q\), define
\begin{align}
 Z_{u,v}(\xi)
 &=
 \sum_{m,n}x_my_n\mathfrak m(m,n)
 \one_{m\equiv-h\overline{\xi}\pmod{u}}
 \one_{n\equiv-h\overline{\xi}\pmod{v}},
 \label{eq:dynamic-pair-fiber}\\
 Z_{q,S}(\xi)
 &=
 \sum_{\substack{u,v\le S\\\operatorname{lcm}(u,v)=q}}
 \vartheta_{R,S}(u)\vartheta_{R,S}(v)Z_{u,v}(\xi).
 \label{eq:dynamic-lcm-fiber}
\end{align}
Inadmissible row--modulus combinations contribute zero.  Let
\begin{equation}
 \overline Z_{u,v}
 =\frac1{\varphi(q)}\sum_{\xi\in\cU_q}Z_{u,v}(\xi),
 \qquad
 D_{u,v}(\xi)=Z_{u,v}(\xi)-\overline Z_{u,v}.
 \label{eq:dynamic-pair-centering}
\end{equation}
Then
\begin{equation}
 \overline Z_{q,S}
 =\frac1{\varphi(q)}\sum_{\xi\in\cU_q}Z_{q,S}(\xi),
 \qquad
 D_{q,S}(\xi)=Z_{q,S}(\xi)-\overline Z_{q,S}.
 \label{eq:dynamic-centering}
\end{equation}
Thus \(D_{q,S}\) is a centered function on \(\cU_q\); the centering is
linear in the divisor-pair decomposition, and explicitly
\begin{equation}
 D_{q,S}
 =
 \sum_{\substack{u,v\le S\\\operatorname{lcm}(u,v)=q}}
 \vartheta_{R,S}(u)\vartheta_{R,S}(v)D_{u,v}.
 \label{eq:dynamic-centered-pair-sum}
\end{equation}

Use
\[
 e(t)=e^{2\pi it},\qquad e_q(t)=e(t/q),\qquad
 c_H(r)=\sum_{\zeta\in(\Z/H\Z)^\times}e_H(r\zeta).
\]
For \(F\in C_c^\infty(\R)\), write
\(\widehat F(t)=\int_\R F(x)e(-tx)\,dx\), and define
\begin{equation}
 \cL_{q,S}
 =\frac{J}{Hq}\sum_{r\ne0}c_H(r)
 \widehat F\left(\frac{rJ}{Hq}\right)
 \sum_{\xi\in\cU_q}Z_{q,S}(\xi)e_q(r\overline H\xi).
 \label{eq:dynamic-scalar-packet}
\end{equation}
Here \(\overline H\) is the inverse of \(H\bmod q\), which exists by
target primitivity.  We write \(\cL_{q,S}^{\rm mean}\) and
\(\cL_{q,S}^{\rm disc}\) for the same linear functional with
\(Z_{q,S}\) replaced by \(\overline Z_{q,S}\) and \(D_{q,S}\),
respectively.  Hence
\[
 \cL_{q,S}=\cL_{q,S}^{\rm mean}+\cL_{q,S}^{\rm disc}.
\]

We retain the TPC-20 hypotheses
\begin{equation}
 J\ge4HX^{\eps_0}
 \label{eq:no-wrap}
\end{equation}
for some fixed \(\eps_0>0\), together with its polynomial bound for the
total absolute Fourier-tail mass.  These assumptions are inherited
interfaces, not consequences of the dynamic cutoff.

Finally, TPC-20's two-scale completion applies simultaneously to all
\(u,v\le S\) under the strict condition
\begin{equation}
 HS\le JX^{-\eta_0}
 \label{eq:S-single-leg}
\end{equation}
for some fixed \(\eta_0>0\).  It removes the two single-indicator and
constant nonzero spectra, leaving the compatible joint packet.  We
refer to \eqref{eq:S-single-leg} as the \emph{Poisson-short condition};
it must not be confused with a restriction to one support quadrant.
